\section{模型评价}
\subsection{问题一模型评价}
双模型框架将热力信息完整时的诊断与仅有雷达时的估计联系起来。模型 a 的参数对应稳定度、谱宽和切变，便于分析各物理过程的作用；模型 b 通过固定参照和时间块验证，量化减少输入后的误差。其主要代价在于参照模型的不确定性会传递到标定结果，因此需要分别报告参照物理检验与统计拟合误差。

在\slot{测试条件}下，模型 b 相对基线的\slot{指标}变化为\slot{数值}；去掉时间特征后变化为\slot{数值}。根据该对照，时间信息对\slot{对应条件}的贡献为\slot{证据支持的解释}。
\subsection{问题二模型评价}
多源融合通过观测算子连接不同采样位置，并通过各向异性尺度区分水平与垂直变化。该框架能够随设备组合调整信息贡献，同时输出观测稀疏区的误差描述。平滑强度过大时会削弱局地峰值，过小时会放大观测噪声，这一取舍可由留站误差与峰值保持情况共同判断。

在设备消融中，去掉\slot{设备}后，\slot{高度或区域}的误差增加\slot{数值}；尺度参数改变\slot{比例}后，关键风险区域的位置变化为\slot{数值}。这些结果用于确定主模型的空间参数和适用范围。
\subsection{问题三模型评价}
模型 d 使用数值模式提供的未来信息，模型 e 使用历史观测的局地演化，两条路线覆盖了题目规定的不同数据条件。路径搜索将风险场转为可重复计算的累计代价，使最终路径能够用独立积分和小图最优值检查。预测时效增加、边界信息不足或风险高值区移动时，路径排序可能发生变化，应结合预报误差与路线敏感性解释结果。

最终结论按三个问题分别落到可交付结果：问题一得到\slot{廓线与误差}，问题二得到\slot{场分布与可靠范围}，问题三得到\slot{预报质量与路径代价}。模型之间的接口、指标单位和时间边界保持一致，支撑程序与结果文件形成逐项对应关系。
\auxfigure{sensitivity}{关键参数敏感性与方法消融比较}{fig:sensitivity}
\begin{table}[H]\centering\auxtablecaption{模型评价与改进方向}
\begin{tabularx}{\linewidth}{lYYY}\toprule
模块 & 已检验内容 & 主要影响因素 & 下一步可验证改进\\\midrule
双模型诊断 & \slot{时间块指标} & 参照误差与信息缺失 & 独立湍流观测标定\\
三维重构 & \slot{留站与消融} & 观测密度与尺度参数 & 增加覆盖或改进误差模型\\
短时预报 & \slot{分时效误差} & 起报信息与局地生消 & 比较更丰富的历史样本\\
路径规划 & \slot{最优值与可行性} & 图离散与风险场误差 & 网格加密与情景重规划\\\bottomrule
\end{tabularx}\end{table}
